\section{Experiment 1: Pull-Request Evaluation}
\label{sec:results-exp1}

This section reports what adding repository context does on 35 pull
requests scored by the rubric. Five Go-only pull requests are excluded
because Joern's Go frontend does not emit the method-level nodes the
queries require. The headline is not that context helps in general (the
structural arm's aggregate effect is small and fragile) but that the two
context types help different things: the structural arm concentrates its advantage in the nine
criteria the mechanism predicts, while retrieval increases rubric coverage
broadly and evenly. That crossing pattern, not any single number, is the
principal result.

\subsection{Paired Comparisons}
\label{sec:results-exp1-paired}

Because every mode reviews every pull request, the informative comparison is
paired. Table~\ref{tab:exp1-deltas} reports each mode against \baseline{}.

\begin{table}[htbp]
    \centering
    \caption{Mean paired differences against \baseline{}, $n = 35$, with
             95\,\% bootstrap intervals.}
    \label{tab:exp1-deltas}
    \small
    \begin{tabular}{lcc}
        \toprule
        Mode & Total /25 $\Delta$ [95\,\% CI] & KG-relevant /9 $\Delta$ [95\,\% CI] \\
        \midrule
        \kg{}     & $+0.71$ [$+0.20$, $+1.26$] & $+0.69$ [$+0.29$, $+1.11$] \\
        \rag{}    & $+1.03$ [$+0.46$, $+1.60$] & $+0.43$ [$-0.03$, $+0.94$] \\
        \hybrid{} & $+0.57$ [$+0.09$, $+1.11$] & $+0.46$ [$+0.03$, $+0.91$] \\
        \bottomrule
    \end{tabular}
\end{table}

The governing pre-specified contrast is the \kg{} gain on the KG-relevant
subscale ($+0.69$, $p = 0.003$, $d_z = +0.56$). \rag{} produces the
largest total-score gain, while \hybrid{} is smaller on both scales. Full
permutation tests and effect sizes are reported in
Appendix~\ref{app:supporting-multiplicity}.
Figure~\ref{fig:exp1-effects} draws both scales together.

\begin{figure}[htbp]
    \centering
    % !TeX root = ../thesis.tex
% GENERATED FILE - do not edit by hand; your changes will be lost.
% Produced by scripts/generate_thesis_figures.py
% Regenerate: python3 scripts/generate_thesis_figures.py
% Computed from:
%   results/JOERN_UNIFIED_FIVE_JUDGE.json (Experiment 1, n = 35)
%
\begin{tikzpicture}
\node[anchor=west,font=\footnotesize] at (-4.5000,0.0000) {\textbf{\kg{}}, KG-relevant /9};
\node[anchor=west,font=\footnotesize] at (-4.5000,-0.5200) {\textbf{\kg{}}, total /25};
\draw[thin,black!55] (-4.5000,-0.7696) -- (6.0000,-0.7696);
\node[anchor=west,font=\footnotesize] at (-4.5000,-1.0400) {\textbf{\rag{}}, KG-relevant /9};
\node[anchor=west,font=\footnotesize] at (-4.5000,-1.5600) {\textbf{\rag{}}, total /25};
\draw[thin,black!55] (-4.5000,-1.8096) -- (6.0000,-1.8096);
\node[anchor=west,font=\footnotesize] at (-4.5000,-2.0800) {\textbf{\hybrid{}}, KG-relevant /9};
\node[anchor=west,font=\footnotesize] at (-4.5000,-2.6000) {\textbf{\hybrid{}}, total /25};
\draw[densely dashed,black!55] (0.3214,0.2340) -- (0.3214,-2.8340);
\draw[kgblue,line width=0.9pt] (1.3571,0.0000) -- (4.2857,0.0000);

\draw[kgblue,line width=0.9pt] (1.3571,-0.0700) -- (1.3571,0.0700);
\draw[kgblue,line width=0.9pt] (4.2857,-0.0700) -- (4.2857,0.0700);
\draw[fill=kgblue,draw=kgblue] (2.7857,0.0000) circle (2.000pt);
\draw[kgblue,line width=0.9pt] (1.0357,-0.5200) -- (4.8214,-0.5200);
\draw[kgblue,line width=0.9pt] (1.0357,-0.5900) -- (1.0357,-0.4500);
\draw[kgblue,line width=0.9pt] (4.8214,-0.5900) -- (4.8214,-0.4500);
\draw[fill=kgblue,draw=kgblue] (2.8571,-0.5200) circle (2.000pt);
\draw[ragorange,line width=0.9pt] (0.2143,-1.0400) -- (3.6785,-1.0400);
\draw[ragorange,line width=0.9pt] (0.2143,-1.1100) -- (0.2143,-0.9700);
\draw[ragorange,line width=0.9pt] (3.6785,-1.1100) -- (3.6785,-0.9700);
\draw[fill=ragorange,draw=ragorange] (1.8571,-1.0400) circle (2.000pt);
\draw[ragorange,line width=0.9pt] (1.9642,-1.5600) -- (6.0000,-1.5600);
\draw[ragorange,line width=0.9pt] (1.9642,-1.6300) -- (1.9642,-1.4900);
\draw[ragorange,line width=0.9pt] (6.0000,-1.6300) -- (6.0000,-1.4900);
\draw[fill=ragorange,draw=ragorange] (4.0000,-1.5600) circle (2.000pt);
\draw[hybridgreen,line width=0.9pt] (0.4285,-2.0800) -- (3.5714,-2.0800);
\draw[hybridgreen,line width=0.9pt] (0.4285,-2.1500) -- (0.4285,-2.0100);
\draw[hybridgreen,line width=0.9pt] (3.5714,-2.1500) -- (3.5714,-2.0100);
\draw[fill=hybridgreen,draw=hybridgreen] (1.9642,-2.0800) circle (2.000pt);
\draw[hybridgreen,line width=0.9pt] (0.6428,-2.6000) -- (4.2857,-2.6000);
\draw[hybridgreen,line width=0.9pt] (0.6428,-2.6700) -- (0.6428,-2.5300);
\draw[hybridgreen,line width=0.9pt] (4.2857,-2.6700) -- (4.2857,-2.5300);
\draw[fill=hybridgreen,draw=hybridgreen] (2.3571,-2.6000) circle (2.000pt);
\draw[thin,black!55] (0.0000,-3.0940) -- (6.5714,-3.0940);
\draw[thin,black!55] (0.3214,-3.0940) -- (0.3214,-3.1740);
\node[anchor=center,font=\scriptsize] at (0.3214,-3.3740) {0};
\draw[thin,black!55] (2.1071,-3.0940) -- (2.1071,-3.1740);
\node[anchor=center,font=\scriptsize] at (2.1071,-3.3740) {0.5};
\draw[thin,black!55] (3.8929,-3.0940) -- (3.8929,-3.1740);
\node[anchor=center,font=\scriptsize] at (3.8929,-3.3740) {1};
\draw[thin,black!55] (5.6786,-3.0940) -- (5.6786,-3.1740);
\node[anchor=center,font=\scriptsize] at (5.6786,-3.3740) {1.5};
\node[anchor=center,font=\footnotesize] at (3.0000,-3.8140) {Paired difference from the diff-only baseline (rubric points)};
\end{tikzpicture}

    \caption{Paired differences against \baseline{} on both scales. Bars are
             95\,\% bootstrap intervals. The modes differ between the total
             rubric and the KG-relevant subscale.}
    \label{fig:exp1-effects}
\end{figure}

The pre-specified \kg{}-versus-\baseline{} contrast remains the governing
analysis. A Friedman sensitivity detects an omnibus mode effect on the total
score ($p = 0.003$) but not on the KG-relevant subscale ($p =
0.084$). The planned contrast remains distinguishable under a
Bonferroni-corrected Wilcoxon test ($p = 0.011$); this is
concordance for that contrast, not evidence of an overall mode effect.
Appendix~\ref{app:supporting-multiplicity} reports the full sensitivity.

\subsection{Criterion-Level Attribution}
\label{sec:results-exp1-attribution}

Table~\ref{tab:exp1-attribution} counts criterion-level wins and losses,
split into the nine KG-relevant criteria and the sixteen others. The sixteen
act as a built-in negative control: a general coverage lift should raise both
groups together.

\begin{table}[htbp]
    \centering
    \caption{Criterion-level wins (W) and losses (L) against \baseline{},
             split into the nine KG-relevant criteria and the sixteen others,
             $n = 35$.}
    \label{tab:exp1-attribution}
    \small
    \begin{tabular}{lcccccc}
        \toprule
        & \multicolumn{3}{c}{KG-relevant (9)} & \multicolumn{3}{c}{Other (16)} \\
        \cmidrule(lr){2-4} \cmidrule(lr){5-7}
        Mode & W & L & net & W & L & net \\
        \midrule
        \kg{}  & 35 & 11 & $+24$ & 14 & 13 & $+1$ \\
        \rag{} & 36 & 21 & $+15$ & 28 & 7 & $+21$ \\
        \bottomrule
    \end{tabular}
\end{table}

This is the chapter's central observation. \kg{} concentrates its net wins
on the criteria it has a mechanism for. Among the other sixteen criteria the
gross movement is not negligible---Q3 gains $+11.4$\,pp while R1 drops $-8.6$\,pp and S1
drops $-5.7$\,pp---but these gains and losses largely cancel, netting to
$+1$. The pattern suggests a displacement effect under the fixed
five-heading format: the review allocates more attention to integration and
triage and less to clarity and input-validation commentary.
The R1 decline is consistent with the human study's later finding that
raters judged the baseline review the better commentary on code clarity
(H5, Section~\ref{sec:results-human}).
\rag{} produces a broad, even lift across both bands. The two modes differ
less in aggregate magnitude than in how selectively they act.

A permutation test confirms the localisation inferentially. If the advantage
were spread at random, the nine KG-relevant criteria would capture 36\,\% of
it. \kg{} places 96\,\% there ($p = 0.003$, $B =
200\,000$); \rag{} places 42\,\% there. This separates the two
mechanisms where aggregate scores cannot.
Figure~\ref{fig:criterion-localisation} shows the per-criterion change in
satisfaction rate for the \kg{} arm relative to the baseline.

\begin{figure}[htbp]
    \centering
    % !TeX root = ../thesis.tex
% GENERATED FILE - do not edit by hand; your changes will be lost.
% Produced by scripts/generate_thesis_figures.py
% Regenerate: python3 scripts/generate_thesis_figures.py
% Computed from:
%   experiments/2026-09-23_five_judge_panel/RESULTS.json (Experiment 1 localisation, n = 40)
%
\begin{tikzpicture}
\node[anchor=east,font=\scriptsize] at (6.5500,0.3060) {$\Delta$\,pp};
\node[anchor=west,font=\footnotesize] at (-5.0000,0.0000) {\emph{KG-relevant band (9 criteria)}};
\node[anchor=west,font=\scriptsize] at (-5.0000,-0.3600) {F3\enspace Integration with existing code};
\fill[kgblue] (1.6710,-0.4824) rectangle (5.2839,-0.2376);
\node[anchor=east,font=\scriptsize] at (6.5500,-0.3600) {+22.9};
\node[anchor=west,font=\scriptsize] at (-5.0000,-0.7200) {F4\enspace Breaking changes or impact};
\fill[kgblue] (1.6710,-0.8424) rectangle (3.9290,-0.5976);
\node[anchor=east,font=\scriptsize] at (6.5500,-0.7200) {+14.3};
\node[anchor=west,font=\scriptsize] at (-5.0000,-1.0800) {T3\enspace Specific test files referenced};
\fill[kgblue] (1.6710,-1.2024) rectangle (3.0258,-0.9576);
\node[anchor=east,font=\scriptsize] at (6.5500,-1.0800) {+8.6};
\node[anchor=west,font=\scriptsize] at (-5.0000,-1.4400) {M3\enspace Public API or interface handling};
\fill[kgblue] (1.6710,-1.5624) rectangle (3.0258,-1.3176);
\node[anchor=east,font=\scriptsize] at (6.5500,-1.4400) {+8.6};
\node[anchor=west,font=\scriptsize] at (-5.0000,-1.8000) {M1\enspace Fit with existing architecture};
\fill[kgblue] (1.6710,-1.9224) rectangle (2.5742,-1.6776);
\node[anchor=east,font=\scriptsize] at (6.5500,-1.8000) {+5.7};
\node[anchor=west,font=\scriptsize] at (-5.0000,-2.1600) {T1\enspace Need for tests discussed};
\fill[kgblue] (1.6710,-2.2824) rectangle (2.1226,-2.0376);
\node[anchor=east,font=\scriptsize] at (6.5500,-2.1600) {+2.9};
\node[anchor=west,font=\scriptsize] at (-5.0000,-2.5200) {T2\enspace Edge cases and error paths in tests};
\fill[kgblue] (1.6710,-2.6424) rectangle (2.1226,-2.3976);
\node[anchor=east,font=\scriptsize] at (6.5500,-2.5200) {+2.9};
\node[anchor=west,font=\scriptsize] at (-5.0000,-2.8800) {C2\enspace Consistency with similar solutions};
\fill[kgblue] (1.6710,-3.0024) rectangle (2.1226,-2.7576);
\node[anchor=east,font=\scriptsize] at (6.5500,-2.8800) {+2.9};
\node[anchor=west,font=\scriptsize,black!55] at (-5.0000,-3.2400) {Q2\enspace Comments anchored to code locations$^{\dagger}$};
\draw[thin,black!55] (-5.0000,-3.5280) -- (6.5500,-3.5280);
\node[anchor=west,font=\footnotesize] at (-5.0000,-3.7260) {\emph{Negative-control band (16 criteria)}};
\node[anchor=west,font=\scriptsize] at (-5.0000,-4.0860) {Q3\enspace Blocking versus non-blocking issues};
\fill[modegrey] (1.6710,-4.2084) rectangle (3.4774,-3.9636);
\node[anchor=east,font=\scriptsize] at (6.5500,-4.0860) {+11.4};
\node[anchor=west,font=\scriptsize] at (-5.0000,-4.4460) {F2\enspace Edge cases in the change identified};
\fill[modegrey] (1.6710,-4.5684) rectangle (2.1226,-4.3236);
\node[anchor=east,font=\scriptsize] at (6.5500,-4.4460) {+2.9};
\node[anchor=west,font=\scriptsize] at (-5.0000,-4.8060) {F1\enspace Change solves the stated problem};
\fill[modegrey] (1.6710,-4.9284) rectangle (2.1226,-4.6836);
\node[anchor=east,font=\scriptsize] at (6.5500,-4.8060) {+2.9};
\node[anchor=west,font=\scriptsize] at (-5.0000,-5.1660) {P1\enspace Obvious inefficiencies};
\fill[modegrey] (1.6710,-5.2884) rectangle (2.1226,-5.0436);
\node[anchor=east,font=\scriptsize] at (6.5500,-5.1660) {+2.9};
\node[anchor=west,font=\scriptsize] at (-5.0000,-5.5260) {S3\enspace Error-handling quality};
\fill[modegrey] (1.6710,-5.6484) rectangle (2.1226,-5.4036);
\node[anchor=east,font=\scriptsize] at (6.5500,-5.5260) {+2.9};
\node[anchor=west,font=\scriptsize,black!55] at (-5.0000,-5.8860) {M2\enspace Pull request too large to review$^{\dagger}$};
\node[anchor=west,font=\scriptsize,black!55] at (-5.0000,-6.2460) {P2\enspace Benchmarks for critical paths$^{\dagger}$};
\node[anchor=west,font=\scriptsize,black!55] at (-5.0000,-6.6060) {Q1\enspace Overall assessment summary$^{\dagger}$};
\node[anchor=west,font=\scriptsize,black!55] at (-5.0000,-6.9660) {Q4\enspace Clarifying questions asked$^{\dagger}$};
\node[anchor=west,font=\scriptsize,black!55] at (-5.0000,-7.3260) {Q5\enspace Rationale given for suggestions$^{\dagger}$};
\node[anchor=west,font=\scriptsize] at (-5.0000,-7.6860) {R2\enspace Unnecessary complexity flagged};
\node[anchor=west,font=\scriptsize,black!55] at (-5.0000,-8.0460) {S2\enspace Hardcoded secrets or credentials$^{\dagger}$};
\node[anchor=west,font=\scriptsize] at (-5.0000,-8.4060) {C1\enspace Project style-guide compliance};
\fill[modegrey!45] (1.2194,-8.5284) rectangle (1.6710,-8.2836);
\node[anchor=east,font=\scriptsize] at (6.5500,-8.4060) {$-$2.9};
\node[anchor=west,font=\scriptsize] at (-5.0000,-8.7660) {R3\enspace Comments explain why, not what};
\fill[modegrey!45] (1.2194,-8.8884) rectangle (1.6710,-8.6436);
\node[anchor=east,font=\scriptsize] at (6.5500,-8.7660) {$-$2.9};
\node[anchor=west,font=\scriptsize] at (-5.0000,-9.1260) {S1\enspace Input validation};
\fill[modegrey!45] (0.7677,-9.2484) rectangle (1.6710,-9.0036);
\node[anchor=east,font=\scriptsize] at (6.5500,-9.1260) {$-$5.7};
\node[anchor=west,font=\scriptsize] at (-5.0000,-9.4860) {R1\enspace Code clarity, naming, function size};
\fill[modegrey!45] (0.3161,-9.6084) rectangle (1.6710,-9.3636);
\node[anchor=east,font=\scriptsize] at (6.5500,-9.4860) {$-$8.6};
\draw[densely dashed,black!55] (1.6710,0.2160) -- (1.6710,-9.6660);
\draw[thin,black!55] (0.0000,-9.8100) -- (5.6000,-9.8100);
\draw[thin,black!55] (0.0903,-9.8100) -- (0.0903,-9.8900);
\node[anchor=center,font=\scriptsize] at (0.0903,-10.0900) {$-$10};
\draw[thin,black!55] (1.6710,-9.8100) -- (1.6710,-9.8900);
\node[anchor=center,font=\scriptsize] at (1.6710,-10.0900) {0};
\draw[thin,black!55] (3.2516,-9.8100) -- (3.2516,-9.8900);
\node[anchor=center,font=\scriptsize] at (3.2516,-10.0900) {10};
\draw[thin,black!55] (4.8323,-9.8100) -- (4.8323,-9.8900);
\node[anchor=center,font=\scriptsize] at (4.8323,-10.0900) {20};
\node[anchor=center,font=\footnotesize] at (2.8000,-10.5300) {Change in satisfaction rate, KG versus baseline (percentage points)};
\end{tikzpicture}

    \caption{Change in criterion satisfaction rate, \kg{} versus \baseline{},
             sorted by magnitude within each band. Daggers mark criteria at
             ceiling or floor in both arms.}
    \label{fig:criterion-localisation}
\end{figure}

\subsection{Inter-Judge Agreement}
\label{sec:results-exp1-agreement}

Pairwise Cohen's $\kappa$ among the five judges ranges from $0.701$ to
$0.893$, substantial to almost perfect under the interpretation in
Section~\ref{sec:bg-stats}. This supports the stability of majority
aggregation, although agreement does not establish that the shared verdict is
correct. Appendix~\ref{app:supporting-agreement} reports all ten pairs,
raw agreement, valid-cell counts, and abstentions.
% Source: experiments/2026-09-23_five_judge_panel/RESULTS.md, lines 32--47.

\subsection{Robustness}
\label{sec:results-exp1-robustness}

The effect was recomputed under several substitutions, varying one element
each time. Figure~\ref{fig:exp1-robustness} draws the estimates on one axis.
The primary row uses the final five-judge panel. The judge rows use one
judge each. The generator and held-out rows are from an earlier
configuration (text-heuristic builder, $n = 40$, three-judge panel) and are
retained as directional evidence. The final row group shows an earlier
generation of the same Joern CPG configuration under the three-judge panel,
which produced $+0.34$ ($p = 0.111$); the difference from the primary row
illustrates the regeneration variance discussed in Section~\ref{sec:repro}. Detailed statistics are in
Appendix~\ref{app:supporting-robustness}.

\begin{figure}[htbp]
    \centering
    % !TeX root = ../thesis.tex
% GENERATED FILE - do not edit by hand; your changes will be lost.
% Produced by scripts/generate_thesis_figures.py
% Regenerate: python3 scripts/generate_thesis_figures.py
% Computed from:
%   experiments/2026-09-23_five_judge_panel/RESULTS.json (headline and individual judges)
%   results/CROSS_GENERATOR_v2.json (generators)
%   results/BOOTSTRAP_STATS_joern_parity.json (builder parity)
%   results/BOOTSTRAP_STATS_confirmatory_clean.json (held-out)
%
\begin{tikzpicture}
\node[anchor=west,font=\footnotesize] at (-5.2000,0.0000) {\emph{Primary configuration}};
\node[anchor=west,font=\footnotesize] at (-5.2000,-0.4600) {Five-judge panel, $n = 35$};
\draw[thin,black!55] (-5.2000,-0.6808) -- (5.4000,-0.6808);
\node[anchor=west,font=\footnotesize] at (-5.2000,-0.9200) {\emph{Single-judge sensitivity}};
\node[anchor=west,font=\footnotesize] at (-5.2000,-1.3800) {gpt-4o only (also the generator)};
\node[anchor=west,font=\footnotesize] at (-5.2000,-1.8400) {gemini-2.5-flash only};
\node[anchor=west,font=\footnotesize] at (-5.2000,-2.3000) {claude-sonnet-4.5 only};
\node[anchor=west,font=\footnotesize] at (-5.2000,-2.7600) {deepseek-v4-pro only};
\node[anchor=west,font=\footnotesize] at (-5.2000,-3.2200) {grok-4.6 only};
\draw[thin,black!55] (-5.2000,-3.4408) -- (5.4000,-3.4408);
\node[anchor=west,font=\footnotesize] at (-5.2000,-3.6800) {\emph{Alternative generators}};
\node[anchor=west,font=\footnotesize] at (-5.2000,-4.1400) {claude-haiku-4.5};
\node[anchor=west,font=\footnotesize] at (-5.2000,-4.6000) {gemini-2.5-flash};
\node[anchor=west,font=\footnotesize] at (-5.2000,-5.0600) {deepseek-chat};
\draw[thin,black!55] (-5.2000,-5.2808) -- (5.4000,-5.2808);
\node[anchor=west,font=\footnotesize] at (-5.2000,-5.5200) {\emph{Regeneration and sample sensitivity}};
\node[anchor=west,font=\footnotesize] at (-5.2000,-5.9800) {Earlier generation, same CPG, three judges};
\node[anchor=west,font=\footnotesize] at (-5.2000,-6.4400) {Held-out pull requests, $n = 12$};
\draw[densely dashed,black!55] (1.1238,-0.2530) -- (1.1238,-6.6470);
\draw[kgblue,line width=0.9pt] (1.7728,-0.4600) -- (5.1107,-0.4600);
\draw[kgblue,line width=0.9pt] (1.7728,-0.5300) -- (1.7728,-0.3900);
\draw[kgblue,line width=0.9pt] (5.1107,-0.5300) -- (5.1107,-0.3900);
\draw[fill=kgblue,draw=kgblue] (3.4418,-0.4600) circle (2.000pt);
\draw[kgblue,line width=0.9pt] (1.9582,-1.3800) -- (5.0180,-1.3800);
\draw[kgblue,line width=0.9pt] (1.9582,-1.4500) -- (1.9582,-1.3100);
\draw[kgblue,line width=0.9pt] (5.0180,-1.4500) -- (5.0180,-1.3100);
\draw[fill=kgblue,draw=kgblue] (3.5345,-1.3800) circle (2.000pt);
\draw[kgblue,line width=0.9pt] (0.2893,-1.8400) -- (4.7398,-1.8400);
\draw[kgblue,line width=0.9pt] (0.2893,-1.9100) -- (0.2893,-1.7700);
\draw[kgblue,line width=0.9pt] (4.7398,-1.9100) -- (4.7398,-1.7700);
\draw[fill=kgblue,draw=kgblue] (2.5146,-1.8400) circle (2.000pt);
\draw[kgblue,line width=0.9pt] (1.4019,-2.3000) -- (3.8126,-2.3000);
\draw[kgblue,line width=0.9pt] (1.4019,-2.3700) -- (1.4019,-2.2300);
\draw[kgblue,line width=0.9pt] (3.8126,-2.3700) -- (3.8126,-2.2300);
\draw[fill=kgblue,draw=kgblue] (2.6073,-2.3000) circle (2.000pt);
\draw[kgblue,line width=0.9pt] (1.8655,-2.7600) -- (4.8326,-2.7600);
\draw[kgblue,line width=0.9pt] (1.8655,-2.8300) -- (1.8655,-2.6900);
\draw[kgblue,line width=0.9pt] (4.8326,-2.8300) -- (4.8326,-2.6900);
\draw[fill=kgblue,draw=kgblue] (3.3490,-2.7600) circle (2.000pt);
\draw[kgblue,line width=0.9pt] (1.4946,-3.2200) -- (4.0908,-3.2200);
\draw[kgblue,line width=0.9pt] (1.4946,-3.2900) -- (1.4946,-3.1500);
\draw[kgblue,line width=0.9pt] (4.0908,-3.2900) -- (4.0908,-3.1500);
\draw[fill=kgblue,draw=kgblue] (2.7927,-3.2200) circle (2.000pt);
\draw[kgblue,line width=0.9pt] (0.4747,-4.1400) -- (2.8854,-4.1400);
\draw[kgblue,line width=0.9pt] (0.4747,-4.2100) -- (0.4747,-4.0700);
\draw[kgblue,line width=0.9pt] (2.8854,-4.2100) -- (2.8854,-4.0700);
\draw[fill=kgblue,draw=kgblue] (1.6801,-4.1400) circle (2.000pt);
\draw[kgblue,line width=0.9pt] (0.7529,-4.6000) -- (5.0180,-4.6000);
\draw[kgblue,line width=0.9pt] (0.7529,-4.6700) -- (0.7529,-4.5300);
\draw[kgblue,line width=0.9pt] (5.0180,-4.6700) -- (5.0180,-4.5300);
\draw[fill=kgblue,draw=kgblue] (2.8854,-4.6000) circle (2.000pt);
\draw[kgblue,line width=0.9pt] (0.4747,-5.0600) -- (3.2563,-5.0600);
\draw[kgblue,line width=0.9pt] (0.4747,-5.1300) -- (0.4747,-4.9900);
\draw[kgblue,line width=0.9pt] (3.2563,-5.1300) -- (3.2563,-4.9900);
\draw[fill=kgblue,draw=kgblue] (1.8655,-5.0600) circle (2.000pt);
\draw[kgblue,line width=0.9pt] (1.0162,-5.9800) -- (3.7718,-5.9800);
\draw[kgblue,line width=0.9pt] (1.0162,-6.0500) -- (1.0162,-5.9100);
\draw[kgblue,line width=0.9pt] (3.7718,-6.0500) -- (3.7718,-5.9100);
\draw[fill=kgblue,draw=kgblue] (2.3959,-5.9800) circle (2.000pt);
\draw[kgblue,line width=0.9pt] (0.5044,-6.4400) -- (4.8326,-6.4400);
\draw[kgblue,line width=0.9pt] (0.5044,-6.5100) -- (0.5044,-6.3700);
\draw[kgblue,line width=0.9pt] (4.8326,-6.5100) -- (4.8326,-6.3700);
\draw[fill=kgblue,draw=kgblue] (2.6703,-6.4400) circle (2.000pt);
\draw[thin,black!55] (0.0000,-6.8770) -- (5.4000,-6.8770);
\draw[thin,black!55] (0.1966,-6.8770) -- (0.1966,-6.9570);
\node[anchor=center,font=\scriptsize] at (0.1966,-7.1570) {$-$0.25};
\draw[thin,black!55] (1.1238,-6.8770) -- (1.1238,-6.9570);
\node[anchor=center,font=\scriptsize] at (1.1238,-7.1570) {0};
\draw[thin,black!55] (2.0510,-6.8770) -- (2.0510,-6.9570);
\node[anchor=center,font=\scriptsize] at (2.0510,-7.1570) {0.25};
\draw[thin,black!55] (2.9782,-6.8770) -- (2.9782,-6.9570);
\node[anchor=center,font=\scriptsize] at (2.9782,-7.1570) {0.5};
\draw[thin,black!55] (3.9054,-6.8770) -- (3.9054,-6.9570);
\node[anchor=center,font=\scriptsize] at (3.9054,-7.1570) {0.75};
\draw[thin,black!55] (4.8326,-6.8770) -- (4.8326,-6.9570);
\node[anchor=center,font=\scriptsize] at (4.8326,-7.1570) {1};
\node[anchor=center,font=\footnotesize] at (2.7000,-7.5970) {Effect on the nine KG-relevant criteria (rubric points)};
\end{tikzpicture}

    \caption{The KG-relevant effect under alternative configurations. Points
             are mean paired effects and bars are 95\,\% bootstrap intervals.}
    \label{fig:exp1-robustness}
\end{figure}

Every recomputation preserves the positive direction, but the magnitude and
uncertainty depend on the judge, generator, and sample. The criterion-level
localisation is more stable than the aggregate magnitude.
